\section{Results and recommendations}

\subsection{Direct answers}

\textbf{Shortening the existing proofs.} Preserve the exact primary K215 and
its public endpoint statements. Adopt the already checked semantic reductions
from the incoming packets, then replace the ordinary source checker's repeated
linear list lookups by the verified balanced representation supplied here.
This changes the certificate implementation without changing the source
program, numerical calculus, binary machine, or consistency statement.
The new theorem is the same all-register source equivalence required by the
existing proof. This is a concrete checked refinement, rather than a proposal
to trust a faster external verifier.

\textbf{An easier additional 215-state proof.} The two-marker family containing
M124 and FM215 is particularly economical when the existing K215 endpoint is
reused. Its entire new mathematical argument is a two-clause tape invariant
and a one-or-three-step local simulation. We package that argument once as
a generic finite checker; three literal instances then occupy 35 lines of
Coq, including an explicit calibration-transfer corollary and assumption
reports. The already supplied CM215 dense/Core bridge may make CM215 the
quickest production integration, even though its marker invariant has one
additional case. A short relative proof does not eliminate the inherited
K215-to-ZFC dependency closure.

\textbf{A cleaner independent enumerator proof.} ZC215's exhaustive raw-code
enumeration admits an especially simple protected-credit proof. We complete
the concrete formalization: the actual Cantor encoding, actual decoder, and
actual increasing candidate order satisfy the required inequality at every
prefix. This is a useful semantic component for a ZC215 formalization; it is
not a new complete ZC215 binary endpoint.

\textbf{214 states.} No valid machine with 214 or fewer working states was
constructed. The expanded local search and its limitations are documented in
Section~\ref{sec:reduction}. The explicit M125 table supplied here has 215
working states and instantiates the provider-125 option already mentioned in
the earlier M124 note. Neither that instance nor the new proof machinery
constitutes a one-state saving.

\subsection{The existing formal baseline}

The accepted native K215 table has 430 rows and SHA-256
\begin{center}\small\ttfamily
9683d6ed9628e6ce407c315b08f5da3d063da94e92706f275ca0f33e8272e4b2.
\end{center}
Its dense form uses working labels 0--214 and halt 215. The maintained
endpoints are:
\begin{quote}\small
Lean: \code{KEndpoint215.blank\_halts\_iff\_zfc\_refutation}\\
Rocq: \code{K215F.CorrectDense215.tm215\_halts\_iff\_ZFC\_inconsistent}
\end{quote}
They express halting iff a finite ZFC refutation exists. Negation gives the
consistency equivalence without taking consistency as a premise
\cite{primary,verification}.

The verification records distinguish the original fresh build from subsequent
incremental simplification. The original Lean run compiled 204 project modules;
the current selected closure has 185. The original Rocq run compiled 237
sources. Phase four rebuilt ten affected Lean modules and reused 175 exactly
bound predecessor objects. Its three focused files decreased from 1,814 to
1,606 lines. Those are existing repository results, not fresh measurements
from this packet \cite{phase4}.

The formal endpoints are external theorems about syntactic ZFC refutations.
The repository explicitly distinguishes them from a proof encoded inside
ZFC's own proof syntax or an internal PA endpoint. The present work preserves
that distinction. Its finite simulations and arithmetic arguments can be
formalized in ZFC at the ordinary mathematical level; it does not deliver a
new internal proof-code object or assert unconditional nonhalting.

\subsection{What the incoming proposals already establish}

All seven archives present at the pinned commit were inspected. Their exact
identities are retained in \code{zero\_credit/incoming\_inventory.json}.
The following reconciliation is essential to avoid repeatedly proposing
work that has already been done \cite{incoming}.

\begin{longtable}{@{}P{.38\textwidth}P{.56\textwidth}@{}}
\toprule
Incoming packet (date suffix omitted) & Existing contribution and boundary\\
\midrule\endhead
\code{K215\_Proof\_Simplification} & Checked CM215 transfer; an exact 1,426-entry
normalized source certificate and source theorem. No fresh whole-endpoint
rebuild.\\[4pt]
\code{K215\_Proof\_Simplification (2)} & Checked CM215 dense transfer; ordinary
zero-offset Lean source lane; compact certificate and Rocq mirror theorem.
Final full-ZFC adapter remains uncompiled in that packet.\\[4pt]
\code{K215\_Proof\_Transport} & Concrete CM215 all-tape and dense-semantics
bridges; a final endpoint wrapper is supplied but uncompiled.\\[4pt]
\code{ZFC\_K215\_Proof\_Transport} & CM215 native transfer, positive-stuttering
theorem, shared-tree and compact-path pilots; representation integration is
separate.\\[4pt]
\code{ZFC215\_Proof\_Simplification} & M124 Lean/Rocq all-tape transfers;
ordinary one-sided source theorem; 88.4601\% certificate-source compression.
Complete M124 representation bridge remains separate.\\[4pt]
\code{zfc-fm215-proof-compression} & FM215 native all-tape Lean transfer and a
234-role decoder certificate. Full endpoint bridge remains separate.\\[4pt]
\code{Structural\_Proof\_Simplification} & Normalized source-graph audit,
semantic seed round-trip proposal, startup sweeps. New formal proposals are
uncompiled; the canonical certificate fails the unchanged raw-pointer checker.\\
\bottomrule
\end{longtable}

The offset symmetry theorem, ordinary zero-offset source proof, compact path
encoding, normalized 1,426 checkpoints, and positive-stuttering transfer are
therefore credited as prior work. Balanced lookup implements an opportunity
left open in the incoming Rocq report. The new marker checker shares the
prior marker argument across a family. The new zero-credit file closes the
concrete encoding and scheduling hypotheses of an earlier abstract prototype.

\subsection{Evidence levels in this packet}

\begin{center}
\begin{tabular}{@{}P{.27\textwidth}P{.32\textwidth}P{.31\textwidth}@{}}
\toprule
Contribution & Fresh validation & Deliberate boundary\\
\midrule
Balanced source theorem & Coq 8.18 compilation, closed assumptions;
recursive check of generic checker & No full 237-source endpoint replay;
concrete theorem checked by compilation\\[4pt]
Marker family & Three modules compiled; six closed reports; recursive
check of concrete instances & Native arbitrary-tape model;
repository representation bridge separate\\[4pt]
Concrete zero credit & 408-line file compiled; nine closed reports;
recursive check passed & Exact arithmetic scheduler, not physical TM execution\\[4pt]
State search & Reproducible C++/Python, controls, exact hashes and counts &
Bounded sufficient proof test; no lower bound\\
\bottomrule
\end{tabular}
\end{center}

All new Coq builds used version 8.18.0 with OCaml 4.14.1. This is different
from the maintained Rocq 9.1.0 lane. The balanced package preserves the
original 9.1 source and its proved compatibility shim. No claim is made that
this task freshly ran either the full Lean endpoint or the full Rocq 9.1
endpoint. Machine-specific proof-assistant results are kept distinct from
Python checks of exact rows and bounded startup runs.
